% Reviewed translation: PDF pages 292--297 / printed pages 278--283.
% Reconstructed and formula-checked against the original scan.
\makeatletter
\gdef\@chapapp{\chaptername}
\gdef\thechapter{\arabic{chapter}}
\gdef\theHchapter{\arabic{chapter}}
\gdef\Hy@chapapp{chapter}
\makeatother
\ctexset{
  chapter/name = {第,章},
  chapter/number = \arabic{chapter},
  chapter/numbering = true
}
\setcounter{chapter}{3}
\chapter{Navier--Stokes 问题的理论与逼近}
\label{chap:navier-stokes-theory-approx}
\counterwithin{equation}{section}
\renewcommand{\thesection}{\arabic{section}}
\renewcommand{\thesubsection}{\thesection.\arabic{subsection}}
\renewcommand{\theequation}{\thesection.\arabic{equation}}
\renewcommand{\thestatement}{\thesection.\arabic{statement}}

\MathBookSourcePage{292}
\section{一类非线性问题}
\label{sec:navier-stokes-nonlinear-problems}
\setcounter{equation}{0}
\setcounter{statement}{0}

本节研究 Chapter~\ref{chap:stokes-foundations} 的
Section~\ref{sec:abstract-variational-problem} 中所分析抽象变分问题的一种非线性推广.
这族非线性问题特别包含 Navier--Stokes 问题.

我们沿用 Section~\ref{sec:abstract-variational-problem} 的记号.
具体地, 考虑两个(实)Hilbert 空间 $X$ 和 $M$,
其范数分别为 $\|\cdot\|_X$ 和 $\|\cdot\|_M$, 以及连续双线性形式
\[
b(\cdot,\cdot):(v,\mu)\in X\times M
\longmapsto b(v,\mu)\in\mathbb R.
\]
非线性通过如下形式引入:
\[
a(\cdot;\cdot,\cdot):(w,u,v)\in X\times X\times X
\longmapsto a(w;u,v)\in\mathbb R,
\]
其中对 $w\in X$, 映射 $(u,v)\mapsto a(w;u,v)$ 是 $X\times X$ 上的连续双线性形式.

于是考虑以下问题, 也称为问题 $(Q)$:

给定 $l\in X'$, 求一对 $(u,\lambda)\in X\times M$, 使
\begin{align}
a(u;u,v)+b(v,\lambda)&=\langle l,v\rangle
&&\forall v\in X,\label{eq:navier-abstract-q-first}\\
b(u,\mu)&=0
&&\forall\mu\in M.\label{eq:navier-abstract-q-second}
\end{align}
引入线性算子 $A(w)\in\mathcal L(X;X')$ ($w\in X$) 和
$B\in\mathcal L(X;M')$, 定义为
\[
\langle A(w)u,v\rangle=a(w;u,v)
\quad\forall u,v\in X,
\]
\[
\langle Bv,\mu\rangle=b(v,\mu)
\quad\forall v\in X,\quad\forall\mu\in M.
\]
在这些记号下, 问题 $(Q)$ 化为: 求 $(u,\lambda)\in X\times M$, 使
\begin{align}
A(u)u+B'\lambda&=l&&\text{于 }X',\label{eq:navier-abstract-q-operator-first}\\
Bu&=0&&\text{于 }M'.\label{eq:navier-abstract-q-operator-second}
\end{align}

\MathBookSourcePage{293}
与线性情形相同, 置 $V=\Ker(B)$, 并将下列称为与问题 $(Q)$ 相联系的问题 $(P)$:

求 $u\in V$, 使
\begin{equation}\label{eq:navier-abstract-p}
a(u;u,v)=\langle l,v\rangle
\qquad\forall v\in V,
\end{equation}
或等价地使
\begin{equation}\label{eq:navier-abstract-p-operator}
\pi A(u)u=\pi l\quad\text{于 }V',
\end{equation}
其中线性算子 $\pi\in\mathcal L(X';V')$ 仍定义为
\[
\langle\pi l,v\rangle=\langle l,v\rangle
\qquad\forall v\in V.
\]

当然, 若 $(u,\lambda)$ 是问题 $(Q)$ 的一个解, 则 $u$ 是问题 $(P)$ 的一个解.
反之, 只要 inf-sup 条件成立, 便可像在线性情形中一样容易地建立逆命题.
因此, 此处真正的困难在于求解非线性问题 $(P)$.
首先需要从下面这个由 Brouwer 给出的经典不动点 Theorem 推出一个简单结论;
我们不给出该 Theorem 的证明.

\begin{Theorem}\label{thm:navier-brouwer-fixed-point}
设 $C$ 是有限维空间的一个非空, 凸且紧的子集,
$F$ 是从 $C$ 到 $C$ 的连续映射. 则 $F$ 至少有一个不动点.
\end{Theorem}

\setcounter{statement}{0}
\begin{Corollary}\label{cor:navier-finite-dimensional-zero}
设 $H$ 是有限维 Hilbert 空间, 其内积记为 $(\cdot,\cdot)$,
相应的范数记为 $|\cdot|$. 设 $\Phi$ 是从 $H$ 到 $H$ 的连续映射,
并具有如下性质: 存在 $\mu>0$, 使
\begin{equation}\label{eq:navier-finite-dimensional-outward-condition}
(\Phi(f),f)\geq0
\qquad\text{对所有满足 }|f|=\mu\text{ 的 }f\in H.
\end{equation}
则存在 $f\in H$, 使
\begin{equation}\label{eq:navier-finite-dimensional-zero}
\Phi(f)=0,\qquad |f|\leq\mu.
\end{equation}
\end{Corollary}

\begin{Proof}
采用反证法. 假设在闭球
$S=\{f\in H;\ |f|\leq\mu\}$ 中 $\Phi(f)\ne0$.
则映射
\[
f\longmapsto-\mu\frac{\Phi(f)}{|\Phi(f)|}
\]
是从 $S$ 到 $S$ 的连续映射. 因为 $H$ 是有限维的,
而 $S$ 显然非空, 凸且紧, 可应用
Theorem~\ref{thm:navier-brouwer-fixed-point}:
存在 $f\in S$, 使
\[
f=-\mu\frac{\Phi(f)}{|\Phi(f)|}.
\]
于是得到一个 $f\in H$, 满足 $|f|=\mu$ 且
\[
(\Phi(f),f)=-\mu|\Phi(f)|<0.
\]
这与 \eqref{eq:navier-finite-dimensional-outward-condition} 矛盾.
\end{Proof}

\MathBookSourcePage{294}
现在可以建立问题 $(P)$ 的如下存在性结果.

\setcounter{statement}{1}
\begin{Theorem}\label{thm:navier-abstract-existence}
假设下列条件成立:

\textup{(i)} 存在常数 $\alpha>0$, 使
\begin{equation}\label{eq:navier-abstract-coercivity}
a(v;v,v)\geq\alpha\|v\|_X^2
\qquad\forall v\in V;
\end{equation}

\textup{(ii)} 空间 $V$ 可分, 且对所有 $v\in V$, 映射
\[
u\longmapsto a(u;u,v)
\]
在 $V$ 上序列弱连续, 即
\begin{equation}\label{eq:navier-abstract-weak-continuity}
\mathop{\mathrm{weak\,lim}}_{m\to\infty}u_m=u\text{ 于 }V
\quad\Longrightarrow\quad
\lim_{m\to\infty}a(u_m;u_m,v)=a(u;u,v)
\quad\forall v\in V.
\end{equation}
则问题 $(P)$ 至少有一个解 $u\in V$.
\end{Theorem}

\begin{Proof}
1) 首先用 Galerkin 方法构造一列近似解.
由于空间 $V$ 可分, 存在 $V$ 中的一列 $(w_m)_{m\geq1}$, 使:

\textup{(i)} 对所有 $m\geq1$, 元素 $w_1,\ldots,w_m$ 线性无关;

\textup{(ii)} $w_i$ 的有限线性组合 $\sum_i v_iw_i$ 在 $V$ 中稠密.

这样的序列 $(w_m)_{m\geq1}$ 称为可分空间 $V$ 的一个``基''.
以 $V_m$ 表示由 $w_1,\ldots,w_m$ 张成的 $V$ 的子空间.
然后用下列问题(称为问题 $(P_m)$)逼近问题 $(P)$:

求 $u_m\in V_m$, 使
\begin{equation}\label{eq:navier-galerkin-problem}
a(u_m;u_m,v)=\langle l,v\rangle
\qquad\forall v\in V_m.
\end{equation}
若置
\[
u_m=\sum_{i=1}^m v_iw_i,
\]
则问题 $(P_m)$ 等价于求解关于 $m$ 个未知量 $v_i$ 的 $m$ 个非线性方程.

现在证明对每个 $m$, 问题 $(P_m)$ 至少有一个解.
引入映射 $\Phi_m:V_m\to V_m$, 定义为
\[
(\Phi_m(v),w_i)=a(v;v,w_i)-\langle l,w_i\rangle,
\qquad1\leq i\leq m,
\]
其中 $(\cdot,\cdot)$ 是 $X$ 中的内积.
因此, $u_m\in V_m$ 是问题 $(P_m)$ 的解, 当且仅当 $\Phi_m(u_m)=0$.
由于
\[
(\Phi_m(v),v)=a(v;v,v)-\langle l,v\rangle
\qquad\forall v\in V_m,
\]
由假设 \eqref{eq:navier-abstract-coercivity} 得
\[
(\Phi_m(v),v)\geq
(\alpha\|v\|_X-\|\pi l\|_{V'})\|v\|_X.
\]

\MathBookSourcePage{295}
因而取 $\mu=\alpha^{-1}\|\pi l\|_{V'}$, 对所有满足
$\|v\|_X=\mu$ 的 $v\in V_m$ 有
\[
(\Phi_m(v),v)\geq0.
\]
此外, 由假设 \eqref{eq:navier-abstract-weak-continuity},
$\Phi_m$ 在 $V_m$ 上连续. 因 $V_m$ 有限维,
可以应用 Corollary~\ref{cor:navier-finite-dimensional-zero}:
问题 $(P_m)$ 至少有一个解 $u_m\in V_m$.
而且, 对所有解 $u_m$,
\[
0=(\Phi_m(u_m),u_m)
\geq(\alpha\|u_m\|_X-\|\pi l\|_{V'})\|u_m\|_X,
\]
所以
\begin{equation}\label{eq:navier-galerkin-uniform-bound}
\|u_m\|_X\leq\frac1\alpha\|\pi l\|_{V'}.
\end{equation}

2) 接着, 对每个 $m$ 任取问题 $(P_m)$ 的一个解,
由此在 $V$ 中构造序列 $(u_m)$. 我们要证明可以抽取一个收敛到问题 $(P)$ 某个解的子列.
由 \eqref{eq:navier-galerkin-uniform-bound}, 序列 $(u_m)$ 在 $V$ 中有界.
因此可抽取一个子列 $(u_{m_p})$, 使
\[
u_{m_p}\rightharpoonup u^*\quad\text{于 }V,
\qquad p\to+\infty.
\]
于是由假设 \eqref{eq:navier-abstract-weak-continuity},
\[
\lim_{p\to\infty}a(u_{m_p};u_{m_p},v)
=a(u^*;u^*,v)
\qquad\forall v\in V.
\]
在 \eqref{eq:navier-galerkin-problem} 中取 $v=w_i$ 和 $m=m_p\geq i$, 得
\[
a(u^*;u^*,w_i)=\langle l,w_i\rangle,
\qquad i\geq1.
\]
因为 $w_i$ 的有限线性组合在 $V$ 中稠密, 得
\[
a(u^*;u^*,v)=\langle l,v\rangle
\qquad\forall v\in V,
\]
故 $u^*$ 是问题 $(P)$ 的一个解.
\end{Proof}

\setcounter{statement}{0}
\begin{Remark}\label{rem:navier-general-existence-context}
显然, Theorem~\ref{thm:navier-abstract-existence} 在更一般的背景下仍成立:
$V$ 是可分 Hilbert 空间; $(u,v)\mapsto a(u,v)$ 是从 $V\times V$ 到 $\mathbb R$ 的映射, 且

\textup{(i)} $a(v,v)\geq\alpha\|v\|_V^2$ ($\forall v\in V$, $\alpha>0$);

\textup{(ii)} 对每个 $u\in V$, $v\mapsto a(u,v)$ 是 $V$ 上的连续线性形式;

\textup{(iii)} 对每个 $v\in V$, $u\mapsto a(u,v)$ 在 $V$ 上序列弱连续.

则给定 $l\in V'$, 至少存在一个 $u\in V$, 使
\[
a(u,v)=\langle l,v\rangle
\qquad\forall v\in V.
\]
\end{Remark}

现在转向问题 $(P)$ 解的唯一性. 这需要比
\eqref{eq:navier-abstract-coercivity} 和
\eqref{eq:navier-abstract-weak-continuity} 更强的假设. 具体假设如下:

\textup{(i)} 双线性形式 $a(w;\cdot,\cdot)$ 关于 $w$ 一致 $V$-椭圆,
即存在常数 $\alpha>0$, 使
\begin{equation}\label{eq:navier-abstract-uniform-ellipticity}
a(w;v,v)\geq\alpha\|v\|_X^2
\qquad\forall v,w\in V;
\end{equation}

\MathBookSourcePage{296}
\textup{(ii)} 映射 $w\mapsto\pi A(w)$ 在 $V$ 中局部 Lipschitz 连续,
即存在连续单调递增函数 $L:\mathbb R_+\to\mathbb R_+$,
使对所有 $\mu>0$,
\begin{equation}\label{eq:navier-abstract-local-lipschitz}
\left|a(w_1;u,v)-a(w_2;u,v)\right|
\leq L(\mu)\|u\|_X\|v\|_X\|w_1-w_2\|_X
\end{equation}
对所有 $u,v\in V$ 及
$w_1,w_2\in S_\mu=\{w\in V;\ \|w\|_X\leq\mu\}$ 成立.

\setcounter{statement}{2}
\begin{Theorem}\label{thm:navier-abstract-uniqueness}
假设 \eqref{eq:navier-abstract-uniform-ellipticity} 和
\eqref{eq:navier-abstract-local-lipschitz} 成立. 则在条件
\begin{equation}\label{eq:navier-abstract-small-data-condition}
\frac{\|\pi l\|_{V'}}{\alpha^2}
L\!\left(\frac{\|\pi l\|_{V'}}{\alpha}\right)<1
\end{equation}
下, 问题 $(P)$ 有唯一解 $u\in V$.
\end{Theorem}

\begin{Proof}
由假设 \eqref{eq:navier-abstract-uniform-ellipticity} 和
Lax--Milgram Theorem~\ref{thm:lax-milgram},
算子 $\pi A(w)\in\mathcal L(V;V')$ 对每个 $w\in V$ 均可逆.
而且 $T(w)=(\pi A(w))^{-1}$ 属于 $\mathcal L(V';V)$ 并满足
\begin{equation}\label{eq:navier-abstract-inverse-bound}
\|T(w)\|_{\mathcal L(V';V)}\leq\frac1\alpha.
\end{equation}
在这些记号下, 问题 $(P)$ 化为
\[
u=T(u)\pi l\quad\text{于 }V.
\]
证明 $v\mapsto T(v)\pi l$ 把 $V$ 映入 $S_\mu$,
并且是 $S_\mu$ 中的严格压缩映射, 其中
$\mu=\alpha^{-1}\|\pi l\|_{V'}$.
对所有 $v\in V$,
\[
\|T(v)\pi l\|_X
\leq\|T(v)\|_{\mathcal L(V';V)}\|\pi l\|_{V'}
\leq\frac1\alpha\|\pi l\|_{V'}=\mu,
\]
故 $T(v)\pi l\in S_\mu$.
接着估计 $u,v\in S_\mu$ 时的 $T(u)-T(v)$.
由恒等式
\[
T(u)-T(v)=T(u)[\pi A(v)-\pi A(u)]T(v)
\]
和 \eqref{eq:navier-abstract-inverse-bound}, 得
\[
\|T(u)-T(v)\|_{\mathcal L(V';V)}
\leq\frac1{\alpha^2}
\|\pi A(v)-\pi A(u)\|_{\mathcal L(V;V')}.
\]
因此 \eqref{eq:navier-abstract-local-lipschitz} 给出
\[
\|(T(u)-T(v))\pi l\|_X
\leq\frac1{\alpha^2}\|\pi l\|_{V'}L(\mu)\|u-v\|_X.
\]
所以, 由 \eqref{eq:navier-abstract-small-data-condition},
映射 $v\mapsto T(v)\pi l$ 是 $S_\mu$ 中的严格压缩映射,
并有唯一不动点 $u\in V$, 即问题 $(P)$ 的唯一解.
\end{Proof}

\setcounter{statement}{1}
\begin{Remark}\label{rem:navier-successive-approximations}
由于 $v\mapsto T(v)\pi l$ 是 $S_\mu$ 中的严格压缩映射,
其不动点 $u$ 可用逐次逼近法计算. 更准确地说,
从任意 $u_0\in S_\mu$ 出发, 构造序列 $(u_m)$:
\[
u_{m+1}=T(u_m)\pi l,
\]
或等价地
\begin{equation}\label{eq:navier-successive-approximation}
a(u_m;u_{m+1},v)=\langle l,v\rangle
\qquad\forall v\in V.
\end{equation}

\MathBookSourcePage{297}
于是
\[
\lim_{m\to\infty}\|u_m-u\|_X=0.
\]
注意不必要求 $u_0\in S_\mu$, 因为对任意 $u_0\in V$,
$u_1=T(u_0)\pi l$ 必然属于 $S_\mu$.
\end{Remark}

最后求解问题 $(Q)$. 如本节开头所述, 其分析与线性情形平行.

\setcounter{statement}{3}
\begin{Theorem}\label{thm:navier-abstract-lagrange-multiplier}
假设双线性形式 $b(\cdot,\cdot)$ 满足 inf-sup 条件
\begin{equation}\label{eq:navier-abstract-inf-sup}
\inf_{\mu\in M}\sup_{v\in X}
\frac{b(v,\mu)}{\|v\|_X\|\mu\|_M}\geq\beta>0.
\end{equation}
则对问题 $(P)$ 的每个解 $u$, 存在唯一 $\lambda\in M$,
使 $(u,\lambda)$ 是问题 $(Q)$ 的一个解.
\end{Theorem}

\begin{Proof}
设 $u\in V$ 是问题 $(P)$ 的一个解. 必须求满足
\eqref{eq:navier-abstract-q-operator-first} 的 $\lambda\in M$.
但 $l-A(u)u$ 属于 $V$ 的极集 $V^0$.
此外, inf-sup 条件 \eqref{eq:navier-abstract-inf-sup} 与
Lemma~\ref{lem:ch4-abstract-div} 蕴含 $B'$ 是从 $M$ 到 $V^0$ 的同构.
因此存在唯一 $\lambda\in M$, 使 $(u,\lambda)$ 是问题 $(Q)$ 的一个解.
\end{Proof}

\setcounter{statement}{2}
\begin{Remark}\label{rem:navier-abstract-mixed-iteration}
现在可扩展 Remark~\ref{rem:navier-successive-approximations} 的迭代格式以求解问题 $(Q)$.
从任意 $u_0\in V$ 出发, 通过求解
\begin{equation}\label{eq:navier-abstract-mixed-iteration}
\left\{
\begin{aligned}
a(u_m;u_{m+1},v)+b(v,\lambda_{m+1})&=\langle l,v\rangle
&&\forall v\in X,\\
b(u_{m+1},\mu)&=0
&&\forall\mu\in M
\end{aligned}
\right.
\end{equation}
构造 $X\times M$ 中的序列 $(u_m,\lambda_m)$.
验证在假设
\eqref{eq:navier-abstract-uniform-ellipticity},
\eqref{eq:navier-abstract-local-lipschitz},
\eqref{eq:navier-abstract-small-data-condition} 和
\eqref{eq:navier-abstract-inf-sup} 下, 对任意 $u_0\in V$ 有
\begin{equation}\label{eq:navier-abstract-mixed-iteration-convergence}
\lim_{m\to\infty}
\left\{\|u_m-u\|_X+\|\lambda_m-\lambda\|_M\right\}=0.
\end{equation}
事实上, $u_m\in V$ 并满足
\eqref{eq:navier-successive-approximation}.
所以由 Remark~\ref{rem:navier-successive-approximations},
\[
\lim_{m\to\infty}\|u_m-u\|_X=0.
\]
另一方面, 利用 \eqref{eq:navier-abstract-q-first} 和
\eqref{eq:navier-abstract-mixed-iteration}, 得
\[
b(v,\lambda_m-\lambda)
=a(u;u,v)-a(u_{m-1};u_m,v)
\qquad\forall v\in X,
\]
并由 \eqref{eq:navier-abstract-inf-sup} 得
\[
\|\lambda_m-\lambda\|_M
\leq\frac1\beta\sup_{v\in X}
\frac{a(u;u,v)-a(u_{m-1};u_m,v)}{\|v\|_X}.
\]
所需结果由 $u_m$ 在 $X$ 中收敛于 $u$ 以及
$a(\cdot;\cdot,\cdot)$ 的连续性性质得到.
\end{Remark}
